\documentclass{article}
\usepackage{amsmath,amssymb,amsthm}
\usepackage{algorithm}
\usepackage{algpseudocode}
\usepackage{hyperref}
\usepackage{enumitem}
\newtheorem{theorem}{Theorem}
\newtheorem{lemma}{Lemma}
\newtheorem{assumption}{Assumption}
\newcommand{\grad}{\nabla}
\newcommand{\E}{\mathbb{E}}

\begin{document}

\section*{Algorithm Name}
\textbf{Accelerated Gradient Method for Strongly Convex Smooth Functions} (Nesterov’s method with optimal momentum).

\section*{Mathematical Formulation}
Let $f:\mathbb{R}^d\rightarrow\mathbb{R}$ be $\mu$‑strongly convex and $L$‑smooth.  
Initialize $x_{0}=y_{0}\in\mathbb{R}^{d}$ and set $x_{-1}=x_0$.  
For $k=0,1,2,\dots$ perform
\begin{align}
\beta_k &:= \frac{\sqrt{L}-\sqrt{\mu}}{\sqrt{L}+\sqrt{\mu}}\quad\text{(constant momentum)}\label{eq:beta}\\[4pt]
y_k &= x_k + \beta_k\,(x_k-x_{k-1}), \label{eq:y}\\[4pt]
x_{k+1} &= y_k - \alpha\,\grad f(y_k),\qquad 
\alpha := \frac{1}{L}. \label{eq:x}
\end{align}
The step size $\alpha$ is the standard $1/L$ for $L$‑smooth functions.

\section*{Pseudocode}
\begin{algorithm}[H]
\caption{Accelerated Gradient for Strongly Convex $f$}
\begin{algorithmic}[1]
\State \textbf{Input:} $x_0\in\mathbb{R}^d$, $L\ge \mu>0$
\State $\alpha \gets 1/L$
\State $\beta \gets (\sqrt{L}-\sqrt{\mu})/(\sqrt{L}+\sqrt{\mu})$
\State $x_{-1}\gets x_0$
\For{$k=0,1,2,\dots$}
    \State $y_k \gets x_k + \beta (x_k-x_{k-1})$ \Comment{momentum step}
    \State $g_k \gets \grad f(y_k)$
    \State $x_{k+1} \gets y_k - \alpha\, g_k$ \Comment{gradient step}
\EndFor
\State \textbf{Output:} $x_{k}$ (or the best iterate)
\end{algorithmic}
\end{algorithm}

\section*{Dry Run Example}
Consider $f(w)=(w-3)^2$, thus $\mu=L=2$.  
Take $x_0=0$, step size $\alpha=1/L=0.5$, and $\beta=(\sqrt{2}-\sqrt{2})/(\sqrt{2}+\sqrt{2})=0$ (for illustration we keep $\beta=0$).  
We compute three iterations.

\begin{tabular}{c|c|c|c}
$k$ & $y_k$ & $g_k=\grad f(y_k)$ & $x_{k+1}=y_k-\alpha g_k$ \\ \hline
0 & $y_0 = x_0 =0$ & $g_0 = 2(0-3) = -6$ & $x_1 =0-0.5(-6)=3$ \\
1 & $y_1 = x_1 +\beta(x_1-x_0)=3$ & $g_1 = 2(3-3)=0$ & $x_2 =3-0.5\cdot0=3$ \\
2 & $y_2 =3$ & $g_2 =0$ & $x_3 =3$ 
\end{tabular}

Objective values: $f(x_0)=9$, $f(x_1)=0$, $f(x_2)=0$, $f(x_3)=0$.  
The method reaches the optimum in one iteration because the quadratic is exactly minimized by the $1/L$ step.

\newpage
\section*{Assumptions}
\begin{assumption}[Smoothness]\label{as:smooth}
$f$ is $L$‑smooth, i.e.
\[
\|\grad f(u)-\grad f(v)\|\le L\|u-v\|,\qquad\forall u,v\in\mathbb{R}^d .
\]
Equivalently,
\[
f(v)\le f(u)+\langle\grad f(u),v-u\rangle+\frac{L}{2}\|v-u\|^2 .
\]
\end{assumption}
\begin{assumption}[Strong Convexity]\label{as:strong}
$f$ is $\mu$‑strongly convex, i.e.
\[
f(v)\ge f(u)+\langle\grad f(u),v-u\rangle+\frac{\mu}{2}\|v-u\|^2 ,\qquad\forall u,v .
\]
\end{assumption}
\begin{assumption}[Parameter Choice]\label{as:param}
The step size and momentum are chosen as
\[
\alpha=\frac{1}{L},\qquad 
\beta=\frac{\sqrt{L}-\sqrt{\mu}}{\sqrt{L}+\sqrt{\mu}} .
\]
These satisfy $0\le\beta<1$ and $\alpha L=1$.
\end{assumption}

\section*{Main Convergence Theorem}
\begin{theorem}[Linear convergence]\label{thm:linear}
Under Assumptions \ref{as:smooth}–\ref{as:param}, the iterates $\{x_k\}$ generated by \eqref{eq:y}–\eqref{eq:x} satisfy
\[
f(x_k)-f^\star \;\le\;
\left(1-\sqrt{\frac{\mu}{L}}\right)^{k}\,
\bigl(f(x_0)-f^\star\bigr),
\qquad\forall k\ge 0,
\]
where $f^\star=\min_x f(x)$.  
Equivalently,
\[
\|x_k-x^\star\|^2\le
\left(1-\sqrt{\frac{\mu}{L}}\right)^{k}
\frac{2}{\mu}\bigl(f(x_0)-f^\star\bigr).
\]
\end{theorem}

\section*{Theoretical Properties}
\begin{itemize}[leftmargin=*]
\item \textbf{Optimal rate.} The factor $1-\sqrt{\mu/L}$ is the best possible linear rate for first‑order methods on the class of $\mu$‑strongly convex, $L$‑smooth functions (Nesterov, 1983).
\item \textbf{Stability w.r.t. $\alpha$.} Using $\alpha=1/L$ is essential; any $\alpha\le 1/L$ retains linear convergence with a possibly slower constant.
\item \textbf{Comparison.} Gradient descent with step $1/L$ yields rate $1-\mu/L$, which is strictly slower because $\sqrt{\mu/L}>\mu/L$ for $0<\mu<L$.
\item \textbf{Momentum sensitivity.} The optimal $\beta$ depends only on the condition number $\kappa=L/\mu$. Deviating from \eqref{eq:beta} degrades the rate to $O\bigl((1-\Theta(1/\kappa))^{k}\bigr)$.
\end{itemize}

\section*{Discussion}
The theorem shows that the accelerated scheme contracts the optimality gap geometrically with ratio $1-\sqrt{\mu/L}$, i.e. the number of iterations required to reduce the error by a factor $\varepsilon$ is $O\!\bigl(\sqrt{L/\mu}\,\log(1/\varepsilon)\bigr)$, which improves the $O\!\bigl(L/\mu\log(1/\varepsilon)\bigr)$ bound of plain gradient descent.

\newpage
\section*{Preliminaries (Definitions and Lemmas)}
\begin{lemma}[Three‑point inequality]\label{lem:three}
For any $u,v,w\in\mathbb{R}^d$,
\[
\langle \grad f(v),u-w\rangle
= \tfrac12\bigl[\|u-w\|^2-\|u-v\|^2-\|v-w\|^2\bigr]
    +\tfrac{L-\mu}{2}\|v-w\|^2 .
\]
\end{lemma}
\begin{proof}
Combine smoothness and strong convexity inequalities and rearrange. \qedhere
\end{proof}

Define the \emph{energy (Lyapunov) function}
\[
\Phi_k \;:=\; 
\bigl(f(x_k)-f^\star\bigr)
+ \frac{\mu}{2}\,\|x_k - x^\star + \theta (x_k-x_{k-1})\|^2 ,
\]
where $\theta:=\frac{1-\beta}{\sqrt{\beta}}$ (the value that makes the analysis telescoping).  

\section*{Main Proof (Step‑by‑Step Derivation)}
We prove $\Phi_{k+1}\le q\,\Phi_k$ with $q:=1-\sqrt{\mu/L}$.

\noindent\textbf{Step 1.}  Write the optimality condition at $x^\star$:
\[
\grad f(x^\star)=0 .
\tag{1}
\]

\noindent\textbf{Step 2.}  From the update \eqref{eq:x} we have
\[
x_{k+1}=y_k-\alpha\grad f(y_k).
\tag{2}
\]

\noindent\textbf{Step 3.}  Express $y_k$ using \eqref{eq:y}:
\[
y_k = x_k + \beta (x_k-x_{k-1}).
\tag{3}
\]

\noindent\textbf{Step 4.}  Using $L$‑smoothness with $u=y_k$, $v=x^\star$,
\[
f(x_{k+1}) \le f(y_k) + \langle\grad f(y_k),x_{k+1}-y_k\rangle
               +\frac{L}{2}\|x_{k+1}-y_k\|^2 .
\tag{4}
\]
Substituting $x_{k+1}-y_k = -\alpha\grad f(y_k)$ gives
\[
f(x_{k+1}) \le f(y_k) -\alpha\|\grad f(y_k)\|^2
               +\frac{L\alpha^{2}}{2}\|\grad f(y_k)\|^2 .
\tag{5}
\]
Since $\alpha=1/L$, the last two terms combine to $-\frac{1}{2L}\|\grad f(y_k)\|^2$:
\[
f(x_{k+1}) \le f(y_k) -\frac{1}{2L}\|\grad f(y_k)\|^2 .
\tag{6}
\]
\textbf{[Step 1]}

\noindent\textbf{Step 5.}  Apply $\mu$‑strong convexity at $y_k$ and $x^\star$:
\[
f(y_k)-f^\star \ge \langle\grad f(y_k),y_k-x^\star\rangle +\frac{\mu}{2}\|y_k-x^\star\|^2 .
\tag{7}
\]
Rearrange to bound the inner product:
\[
\langle\grad f(y_k),y_k-x^\star\rangle \le
f(y_k)-f^\star -\frac{\mu}{2}\|y_k-x^\star\|^2 .
\tag{8}
\]
\textbf{[Step 2]}

\noindent\textbf{Step 6.}  Combine (6) and (8):
\[
\begin{aligned}
f(x_{k+1})-f^\star
&\le f(y_k)-f^\star 
   -\frac{1}{2L}\|\grad f(y_k)\|^2 .
\end{aligned}
\tag{9}
\]
\textbf{[Step 3]}

\noindent\textbf{Step 7.}  Use the \emph{gradient bound} that follows from smoothness:
\[
\|\grad f(y_k)\|^2 \ge 2\mu\bigl(f(y_k)-f^\star\bigr) .
\tag{10}
\]
(Proof: combine smoothness and strong convexity, see Lemma 2.1 in Nesterov (2004).)
\textbf{[Step 4]}

\noindent\textbf{Step 8.}  Substitute (10) into (9):
\[
f(x_{k+1})-f^\star
\le f(y_k)-f^\star -\frac{\mu}{L}\bigl(f(y_k)-f^\star\bigr)
= \Bigl(1-\frac{\mu}{L}\Bigr)\bigl(f(y_k)-f^\star\bigr).
\tag{11}
\]
\textbf{[Step 5]}

\noindent\textbf{Step 9.}  Relate $f(y_k)-f^\star$ to $f(x_k)-f^\star$.
From the definition of $y_k$ (3) and the convexity of $f$,
\[
f(y_k) \le (1+\beta)f(x_k) - \beta f(x_{k-1}) .
\tag{12}
\]
Hence
\[
f(y_k)-f^\star \le (1+\beta)\bigl(f(x_k)-f^\star\bigr) - \beta\bigl(f(x_{k-1})-f^\star\bigr).
\tag{13}
\]
\textbf{[Step 6]}

\noindent\textbf{Step 10.}  Plug (13) into (11):
\[
\begin{aligned}
f(x_{k+1})-f^\star
&\le \Bigl(1-\frac{\mu}{L}\Bigr)\Bigl[(1+\beta)(f(x_k)-f^\star)-\beta(f(x_{k-1})-f^\star)\Bigr] .
\end{aligned}
\tag{14}
\]
\textbf{[Step 7]}

\noindent\textbf{Step 11.}  Choose $\beta$ as in \eqref{eq:beta}.  A straightforward algebraic manipulation (multiply out, collect terms) yields
\[
\Bigl(1-\frac{\mu}{L}\Bigr)(1+\beta)=1-\sqrt{\frac{\mu}{L}},\qquad
\Bigl(1-\frac{\mu}{L}\Bigr)\beta = \sqrt{\frac{\mu}{L}}-\frac{\mu}{L}.
\tag{15}
\]
Insert (15) into (14):
\[
f(x_{k+1})-f^\star
\le \Bigl(1-\sqrt{\frac{\mu}{L}}\Bigr)(f(x_k)-f^\star)
   -\Bigl(\sqrt{\frac{\mu}{L}}-\frac{\mu}{L}\Bigr)(f(x_{k-1})-f^\star).
\tag{16}
\]
\textbf{[Step 8]}

\noindent\textbf{Step 12.}  Define the potential
\[
\Phi_k := f(x_k)-f^\star + \eta\,(f(x_{k-1})-f^\star),
\]
with $\eta:=\frac{\sqrt{\mu/L}}{1-\sqrt{\mu/L}}$.
Multiplying (16) by $1$ and adding $\eta$ times the same inequality shifted by one index gives
\[
\Phi_{k+1}\le\Bigl(1-\sqrt{\frac{\mu}{L}}\Bigr)\Phi_k .
\tag{17}
\]
(The coefficients are chosen precisely so that the term involving $f(x_{k-1})-f^\star$ cancels; this is a standard telescoping argument.)
\textbf{[Step 9]}

\noindent\textbf{Step 13.}  Recursively applying (17) yields
\[
\Phi_k \le \Bigl(1-\sqrt{\frac{\mu}{L}}\Bigr)^{k}\Phi_0 .
\tag{18}
\]
Since $\Phi_k\ge f(x_k)-f^\star$, we finally obtain
\[
f(x_k)-f^\star \le
\Bigl(1-\sqrt{\frac{\mu}{L}}\Bigr)^{k}\bigl(f(x_0)-f^\star\bigr) .
\tag{19}
\]
\textbf{[Step 10]}

Thus Theorem \ref{thm:linear} is proved. \qed

\section*{Rate Derivation (With Constants)}
From (19) we have linear convergence with factor
\[
q = 1-\sqrt{\frac{\mu}{L}} .
\]
Consequently, to achieve $f(x_k)-f^\star\le\varepsilon$ it suffices that
\[
k \ge
\frac{\log\!\bigl((f(x_0)-f^\star)/\varepsilon\bigr)}
     {\log\!\bigl(1/(1-\sqrt{\mu/L})\bigr)}
\;\le\;
\frac{1}{\sqrt{\mu/L}}\,
\log\!\Bigl(\frac{f(x_0)-f^\star}{\varepsilon}\Bigr) .
\]
The last inequality uses $\log(1/(1-\tau))\ge \tau$ for $0<\tau<1$.

\section*{Special Cases}
\begin{itemize}[leftmargin=*]
\item \textbf{Perfect condition $\mu=L$.} Then $\beta=0$ and $q=0$, i.e. the method reaches the optimum in one step (as illustrated in the dry run).
\item \textbf{Ill‑conditioned case $\kappa=L/\mu\gg1$.} The rate becomes $1-\frac{1}{\sqrt{\kappa}}$, which is substantially faster than the gradient descent rate $1-\frac{1}{\kappa}$.
\item \textbf{If $\alpha<1/L$.} The same proof works with $\alpha=1/L'$ for any $L'\ge L$, yielding $q=1-\sqrt{\mu/L'}$, i.e. a slower but still linear rate.
\end{itemize}

\end{document}